\section{Discussion, Limitations \& Conclusion}
\label{sec:discussion}

\subsection{Discussion of Findings}
% TODO (Prompt 2): Comprehensive analysis of full training results once executed on Kaggle GPU.
% Topics to expand:
% 1. Comparative analysis of Gradient Boosted Trees (LightGBM) vs Deep 1D-ResNet representation learning.
% 2. Error mode breakdown: high-frequency attenuation smearing in long lines and edge overlapping in near-load regimes.
% 3. Effectiveness of the Savitzky-Golay FIR derivative channel in regularizing edge localization.
% 4. Interpretation of physics-masked multi-task loss dynamics and gradient stability across conflicting targets.
% 5. Explainability insights from SHAP feature attributions on impedance discontinuities.

The benchmark results validate the necessity of physics-informed representations when diagnosing distributed transmission line systems. While classical edge-detection mechanisms break down under dispersion and multiple overlapping echoes, machine learning models exhibit strong resilience by capturing holistic wave dynamics across the entire normalized propagation window.

\subsection{Limitations and Future Scope}
% TODO (Prompt 2): In-depth discussion of physical limitations and future research avenues:
% 1. Extension from single-discontinuity cables to branched wiring topologies and complex network meshes (e.g., aircraft harnesses).
% 2. Incorporation of variable temperature profiles along the line length altering local resistance and dielectric permittivity.
% 3. Adaptation to non-TEM waveguides, twisted pairs, and multi-conductor transmission lines with modal dispersion.
% 4. Few-shot domain adaptation to close the sim-to-real gap on uncalibrated commercial TDR instruments.

\subsection{Conclusion}
% TODO (Prompt 2): Final concluding summary emphasizing key quantitative milestones from the full dataset run.

In this work, we presented a complete end-to-end framework for machine-learning-based transmission line fault and load identification via time-domain reflectometry. Grounded in an exact frequency-domain telegrapher cascade simulator capturing copper skin effect and dielectric losses, our 3-channel standardized representation enables simultaneous fault classification (T1), precise metric localization (T2), and complex reactive termination estimation (T3). The physics-masked multi-task loss guarantees uncorrupted training gradients, while empirical benchmarking confirms robust generalization across severe noise levels, out-of-distribution physical shifts, and real-world laboratory coaxial cable measurements.

